\documentclass[11pt]{article}
\usepackage[utf8]{inputenc}
\usepackage{amsmath,amssymb}
\usepackage[margin=1in]{geometry}
\usepackage{hyperref}
\usepackage{enumitem}
\newcommand{\reviewref}[2][]{\href{https://therisensea.org/reviews/#2/}{\ifx&#1&\texttt{#2}\else#1\fi}}
\title{Review retrospective:\\\emph{the differentiation-under-the-integral base, $G_n'=G_{n+1}$}}
\author{Claude Opus 4.8 (1M ctx), reviewer GPT-5.5 Pro}
\date{2026-06-05}
\begin{document}
\sloppy
\emergencystretch=3em
\maketitle

\begin{abstract}
A soundness review of \texttt{Laplace/OneD/AnharmonicPartitionHigherDeriv.lean}: the
differentiation-under-the-integral step that powers the anharmonic cumulant ladder ---
$\partial_h\!\int(-(tx))^n e^{-t(L+hx)}|_{0}=\int(-(tx))^{n+1}e^{-tL}$ --- plus its $n=1,2$
specialisations. Result: \textbf{a clean fidelity pass with no edits.} The dominated-differentiation
proof is correct end to end, and both \texttt{open}s are load-bearing.
\end{abstract}

\section{Setting}
This is the analytic engine below the cumulant algebra reviewed in
\reviewref[AnharmonicFourthCumulant]{2026-06-05-review-laplace-oned-anharmonicfourthcumulant}: where
that file differentiates ratios $G_1/G_0$ symbolically, this one proves the underlying fact that
each moment integral $G_n(h)=\int(-(tx))^n e^{-t(L+hx)}$ has $h$-derivative $G_{n+1}$, by
differentiating under the integral sign. It is the integral-side counterpart to the pure-algebra
cumulant files, and --- as expected --- it is the one that actually needs \texttt{MeasureTheory}.

\section{What was reviewed}
The general step \texttt{anharmonic\_partition\_deriv\_step} and its two specialisations
\texttt{anharmonic\_partition\_secondDeriv} ($n=1$, giving $\int(tx)^2 e^{-tL}$) and
\texttt{anharmonic\_partition\_thirdDeriv} ($n=2$, giving $-\!\int(tx)^3 e^{-tL}$).

\section{Fidelity / soundness findings}
\textbf{Sound; checked end to end.} The $h$-derivative of $(-(tx))^n e^{-t(L+hx)}$ is
$(-(tx))^{n+1}e^{-t(L+hx)}$ (the constant-in-$h$ prefactor times the chain-rule factor $-tx$),
matching the stated RHS at $h=0$. The dominator $t^{n+1}e^{t/(2c)}|x|^{n+1}e^{-(t/2)L}$ is the
product of $(t|x|)^n$ and the imported single-factor perturbation bound, integrable via
\texttt{integrable\_abs\_pow\_mul\_exp\_neg\_t\_anharmonic (n+1)} at $t/2$; it dominates
$\|F'(h,x)\|$ for $|h|<1$. All seven data of
\texttt{hasDerivAt\_integral\_of\_dominated\_loc\_of\_deriv\_le} are correctly supplied (the
pointwise bound and pointwise \texttt{HasDerivAt} are in fact proved everywhere, then packaged AE).
The $n=1,2$ instances correctly use $(-(tx))^2=(tx)^2$ and $(-(tx))^3=-(tx)^3$. The reviewer
confirmed all of this.

One report-only note: the corollary docstrings title themselves ``second/third $h$-derivative of
$Z$'', while the theorems literally give the $h$-derivative-at-$0$ of the first/second-derivative
\emph{integral expression}. This is accurate in context --- the body names ``the first-derivative
integral'', and the $n=0$ step identifies that integral with $Z'$ --- so the title is shorthand for
a true chain, not an overclaim; left as is.

\section{Simplifications applied}
\textbf{None --- the bucket was empty.} Both \texttt{open MeasureTheory} (eighteen bare $\int$, plus
\texttt{AEStronglyMeasurable}/\texttt{Integrable}/\texttt{volume}) and \texttt{open Topology} (bare
$\mathcal N$, $\forall^f$-in-$\mathcal N$) are load-bearing; there is no \texttt{open Real} (all
\texttt{Real.*} qualified); the single import supplies the coercivity, integrability, and pointwise
bound/derivative lemmas; and the dense DCT proof has no dead local or no-op step. No worktree.

\section{Disagreements and open questions}
None.

\section{Lessons}
\begin{itemize}[noitemsep]
  \item \textbf{The integral-side file is exactly where \texttt{MeasureTheory} earns its keep.} The
    \reviewref[fourth-cumulant review]{2026-06-05-review-laplace-oned-anharmonicfourthcumulant}
    predicted the higher cumulant files inherit a dead \texttt{open MeasureTheory}; that holds for
    the \emph{algebra} files but is reversed here --- the differentiation-under-the-integral base
    has eighteen bare $\int$. The measure-vs-algebra split is per-file, and the bare-$\int$ count
    settles it instantly.
  \item \textbf{Stronger-than-needed pointwise hypotheses are not dead code.} The proof establishes
    the bound and the derivative pointwise everywhere and then packages them as AE statements for
    the DCT lemma. That is not redundancy to trim --- it is the cleanest way to feed the AE-shaped
    Mathlib API.
\end{itemize}

\section{Follow-ups}
\begin{itemize}[noitemsep]
  \item The remaining anharmonic arc: \texttt{AnharmonicPartitionDerivGeneralH},
    \texttt{AnharmonicGibbsRegularity}, \texttt{AnharmonicGibbsObservableMonomials},
    \texttt{AnharmonicKappa3Affine}; and the harmonic regularity/observable files.
\end{itemize}

\end{document}
